\documentclass[a4paper,12pt]{article}

\usepackage{cite}
\usepackage{amsmath,amssymb,amsfonts}
\usepackage{algorithmic}
\usepackage{graphicx}
\usepackage{textcomp}
\usepackage{xcolor}
\usepackage{booktabs}
\usepackage{url}
\def\BibTeX{{\rm B\kern-.05em{\sc i\kern-.025em b}\kern-.08em
    T\kern-.1667em\lower.7ex\hbox{E}\kern-.125emX}}
\begin{document}

\title{Assignment 1 : Supervised Machine Learning %\\
%{\footnotesize \textsuperscript{*}Note: Sub-titles are not captured for https://ieeexplore.ieee.org  and
%should not be used}
%\thanks{Identify applicable funding agency here. If none, delete this.}
}

\author{Renjiro Owen\\
\textit{Victoria University Wellington}\\
Wellington, New Zealand\\
owenrenj@myvuw.ac.nz}

\maketitle

\section{Decision Tree Method}

\subsection{Results}

The decision tree algorithm was trained using the rtg A, rtg B, and rtg C training dataset. 
The generated decision tree was constructed using entropy-based Information Gain, 
with a stopping threshold of $IG < 0.00001$.
The overall training classification accuracy achieved was:

\begin{equation}
	Accuracy_{rtgA} = \text{100}\%
\end{equation}

\begin{equation}
	Accuracy_{rtgB} = \text{100}\%
\end{equation}

\begin{equation}
	Accuracy_{rtgC} = \text{100}\%
\end{equation}

The classification accuracy was calculated using:

\begin{equation}
	Accuracy = \frac{\text{Number of correctly classified instances}}
	{\text{Total number of instances}}
\end{equation}

\subsection{Minimal Entropy}

For each dataset, the feature resulting in the minimum entropy was identified by
calculating the weighted entropy after splitting the entire training dataset on
each available feature. The feature with the lowest resulting entropy represents
the split that produces the purest separation of class labels and therefore
provides the greatest reduction in uncertainty.

For \textit{rtg A}, Feature 0 produced the minimum entropy value of 0.285506 with
an Information Gain of 0.116673. Therefore, Feature 0 was selected as the root
splitting feature of the generated decision tree. For \textit{rtg B}, Feature 3
resulted in the lowest entropy value of 0.900495 with an Information Gain of
0.053199, making it the most effective feature for the initial split. For
\textit{rtg C}, Feature 3 achieved an entropy value of 0.000000 and an Information
Gain of 0.904381, indicating that this feature completely separates the classes
in the training dataset.

These observations were made by comparing the entropy values obtained from all
features. The feature with the smallest entropy (or equivalently the highest
Information Gain) was considered the most informative feature and was chosen by
the decision tree algorithm for splitting the dataset.

\subsection{Unused Feature}
The features utilized by the decision tree were different for each dataset. For
\textit{rtg A}, Feature 3 was not used in constructing the decision tree. This
occurred because the other features provided sufficient Information Gain to
separate the classes, allowing the tree to reach pure leaf nodes before Feature
3 was considered. Since the stopping criteria were satisfied once the nodes
became sufficiently pure, Feature 3 did not contribute additional useful
information for classification.

For \textit{rtg B}, all available features were utilized in the construction of
the decision tree. This indicates that each feature contributed useful
information during the splitting process. The tree required all features to
further separate the data until the stopping criteria were reached.

For \textit{rtg C}, only Feature 3 was used in constructing the decision tree,
while the remaining features were excluded. This was because Feature 3 provided
a very strong separation of the classes, achieving an entropy of 0 and a high
Information Gain of 0.904381. After splitting on Feature 3, the resulting nodes
were already pure, meaning additional features were unnecessary. Therefore, the
remaining features were not selected because they did not provide any further
improvement to the classification.

Overall, features were excluded from the tree-building process when they did not
provide enough additional Information Gain or when previous splits had already
produced pure leaf nodes. The decision tree only selects features that reduce
class uncertainty and ignores redundant or unnecessary attributes.

\subsection{Rtg C Discussion}
\subsubsection{Generalisation Ability of the Generated Tree}

The generated decision tree is not expected to perform well on unseen data. The
tree structure indicates that Feature 3 alone is able to completely separate the
training instances, resulting in a very simple tree where the split is entirely
dependent on a single feature. Although this produces perfect classification
performance on the training dataset, it may not generalise well to unseen data
because the model relies heavily on Feature 3.

If new instances contain feature values that were not represented in the training
dataset, the tree may not have sufficient information from other features to make
a reliable prediction. Therefore, the tree structure suggests a risk of poor
generalisation due to its dependency on a single feature rather than learning
patterns from multiple attributes.


\subsubsection{Unused Features in Tree Construction}

For the rtg C dataset, only Feature 3 was utilized in constructing the decision
tree, while the remaining features were not selected. This occurred because
Feature 3 provided a very strong separation between the two classes, achieving an
entropy of 0 and a high Information Gain of 0.904381.

Since splitting on Feature 3 resulted in pure leaf nodes, there was no need for
the decision tree to consider additional features. The other attributes may
contain useful information; however, their contribution was unnecessary after
Feature 3 had already completely separated the training data. Therefore, these
features were excluded because they did not provide additional Information Gain
during the tree-building process.


\subsubsection{Potential Preprocessing Improvements}

Although a test dataset is not available, several preprocessing approaches could
be considered to improve the generalisation ability of the generated tree.
Firstly, Feature 3 could be investigated further because the tree's complete
dependency on this single attribute may indicate that the dataset contains a
dominant or potentially artificial feature. Removing Feature 3 before training
could encourage the model to learn from other attributes and produce a more
generalised decision boundary.

Additionally, feature selection techniques could be applied to identify and
retain informative attributes while reducing dependency on a single feature.
Tree pruning or increasing the Information Gain stopping threshold could also be
considered to prevent the tree from creating overly specific splits based only on
the training data.

These preprocessing steps may reduce overfitting and improve the ability of the
decision tree to classify unseen instances.

\section{Exploring Classification Algorithms}
\subsection{Missing Data Processing and Data Normalization}
For the missing-data handling experiment, two imputation strategies, mean imputation and median imputation, were applied to the dataset before training the Decision Tree classifier. The target variable was set as diagnosis, and the dataset was split into training and testing sets using a 70:30 ratio with a fixed random state of 42 to ensure reproducibility. A Decision Tree classifier with a single plausible hyperparameter configuration (maximum depth = 8) was used to evaluate the effect of each imputation method. The results showed that both mean and median imputation achieved the same classification accuracy of 91.23\%, indicating that the choice of imputation strategy had no significant impact on the model performance for this dataset. Since mean imputation achieved the highest accuracy (91.23\%), it was selected as the best missing-data handling method. The completed dataset generated using the selected imputation approach was then used for all subsequent tasks.

\begin{table}[h]
	\centering
	\begin{tabular}{lc}
		\hline
		\textbf{Imputation Method} & \textbf{Accuracy (\%)} \\
		\hline
		Mean Imputation & 91.23 \\
		Median Imputation & 91.23 \\
		\hline
	\end{tabular}
	\caption{Decision Tree Accuracy Comparison Using Different Missing-Data Imputation Strategies}
	\label{tab:decision_tree_imputation}
\end{table}

Standardization and Min-Max scaling were applied to the completed dataset to investigate the impact of feature normalization on KNN and Decision Tree classifiers. The dataset was evaluated using a fixed train-test split with one plausible hyperparameter configuration for each model: K = 5 for KNN and maximum depth = 8 for the Decision Tree. The results are presented in Table \ref{tab:normalization_comparison}. Without normalization, the KNN classifier achieved an accuracy of only 69.59\%, while applying standardization and Min-Max scaling significantly improved its performance to 96.49\% and 97.08\%, respectively. This demonstrates that KNN is highly affected by feature scaling because it relies on distance calculations between samples. In contrast, the Decision Tree classifier achieved the same accuracy of 91.23\% across all three datasets, indicating that normalization had no noticeable effect on its performance. This is expected because Decision Trees make decisions based on feature thresholds rather than distance measurements. Therefore, normalization techniques have a much greater impact on KNN compared with Decision Tree models, with Min-Max scaling providing the best performance improvement for KNN in this experiment.

\begin{table}[h]
	\centering
	\setlength{\tabcolsep}{4pt}
	\begin{tabular}{lcc}
		\hline
		\textbf{Normalization Method} & \textbf{KNN (\%)} & \textbf{Decision Tree (\%)} \\
		\hline
		No Normalization & 69.59 & 91.23 \\
		Standardization & 96.49 & 91.23 \\
		Min-Max Scaling & 97.08 & 91.23 \\
		\hline
	\end{tabular}
	\caption{Accuracy Comparison of KNN and Decision Tree Models with Different Normalization Techniques}
	\label{tab:normalization_comparison}
\end{table}

\subsection{Classifier Exploration and Hyperparameter Tuning}
The performance of four classification algorithms, KNN, Decision Tree, AdaBoost, and Random Forest, was evaluated by varying their key hyperparameters while keeping the dataset split consistent. The results are shown in Table~\ref{tab:classifier_hyperparameter_comparison}. For KNN, changing the number of neighbors had a noticeable impact on performance. The best result was achieved with n\_neighbors = 9, producing the highest accuracy of 97.66\% and an F1 score of 0.9677. Smaller values of $k$ can make the model more sensitive to noise, while larger values may oversmooth the decision boundary and reduce classification performance.

For the Decision Tree classifier, increasing the maximum tree depth from 2 to 8 improved accuracy from 90.64\% to 91.23\%. However, further increasing the depth to 14 did not provide additional improvement, suggesting that the model had already captured the important patterns in the dataset and deeper trees did not contribute to better generalization. The similar precision, recall, and F1 scores between depth values of 8 and 14 indicate that additional complexity did not improve predictive performance.

AdaBoost showed continuous improvement as the number of weak learners increased. Increasing n\_estimators from 10 to 30 improved accuracy from 94.74\% to 96.49\%, with corresponding improvements in recall and F1 score. This occurs because boosting combines multiple weak learners to reduce prediction errors and create a stronger ensemble model.

Random Forest also benefited from increasing the number of trees. The best performance was achieved with n\_estimators = 50, reaching an accuracy of 97.08\%, precision of 1.0000, recall of 0.9219, and an F1 score of 0.9594. However, increasing the number of trees further to 60 slightly reduced performance, indicating that additional estimators provided limited benefit.

Overall, KNN with n\_neighbors = 9 achieved the highest accuracy of 97.66\%, making it the best-performing classifier in this experiment. Random Forest with 50 trees also achieved strong performance, demonstrating the effectiveness of ensemble learning methods. The superior performance of KNN is likely due to the normalized feature values and the dataset's clear separation between classes, allowing distance-based classification to identify similar samples effectively. Random Forest performed well because combining multiple decision trees reduces overfitting and improves model stability compared with a single Decision Tree. AdaBoost also achieved competitive results by iteratively improving weak learners, while the standalone Decision Tree showed lower performance due to its higher tendency to overfit and sensitivity to the selected tree depth.
\begin{table}[h]
	\centering
	\footnotesize
	\setlength{\tabcolsep}{3pt}
	\begin{tabular}{llccccc}
		\hline
		\textbf{Classifier} & \textbf{Hyperparameter} & \textbf{Value} & \textbf{Accuracy (\%)} & \textbf{Precision} & \textbf{Recall} & \textbf{F1 Score} \\
		\hline
		KNN & n\_neighbors & 3  & 97.08 & 1.0000 & 0.9219 & 0.9594 \\
		 & & 9  & 97.66 & 1.0000 & 0.9375 & 0.9677 \\
		 & & 15 & 95.91 & 1.0000 & 0.8906 & 0.9421 \\
		 & & 21 & 96.49 & 1.0000 & 0.9063 & 0.9508 \\
		\hline
		Tree & max\_depth & 2  & 90.64 & 0.9444 & 0.7969 & 0.8644 \\
		 & & 8  & 91.23 & 0.9298 & 0.8281 & 0.8760 \\
		 & & 14 & 91.23 & 0.9298 & 0.8281 & 0.8760 \\
		\hline
		AdaBoost & n\_estimators & 10 & 94.74 & 0.9825 & 0.8750 & 0.9256 \\
		 & & 20 & 95.91 & 0.9831 & 0.9063 & 0.9431 \\
		 & & 30 & 96.49 & 0.9833 & 0.9219 & 0.9516 \\
		\hline
		Forest & n\_estimators & 10 & 95.32 & 0.9828 & 0.8906 & 0.9344 \\
		 & & 30 & 95.91 & 0.9831 & 0.9063 & 0.9431 \\
		 & & 50 & 97.08 & 1.0000 & 0.9219 & 0.9594 \\
		 & & 60 & 96.49 & 1.0000 & 0.9063 & 0.9508 \\
		\hline
	\end{tabular}
	\caption{Classifier Performance with Different Hyperparameter Configurations}
	\label{tab:classifier_hyperparameter_comparison}
\end{table}
\subsection{Feature Engineering and Dimensionality Reduction}
\subsubsection{Feature Selection using Pearson Correlation}

Pearson correlation was applied to identify the relationship between each feature and the target variable (diagnosis). A correlation threshold of 0.6 was selected, where only features with an absolute correlation coefficient greater than 0.6 were retained. Features satisfying this condition were considered highly relevant for classification. The Decision Tree classifier was then trained using the reduced feature set, and its performance was compared with the model trained using all original features.

The correlation coefficients between the features and the diagnosis variable are shown in Table~\ref{tab:correlation}. Based on the threshold value of 0.6, ten features were selected: radius\_mean, perimeter\_mean, area\_mean, concavity\_mean, concave points\_mean, radius\_worst, perimeter\_worst, area\_worst, concavity\_worst, and concave points\_worst. These selected features represent measurements that have the strongest relationship with the diagnosis outcome.

\begin{table}
	\centering
	\begin{tabular}{l r}
		\hline
		\textbf{Feature} & \textbf{Correlation Coefficient} \\
		\hline
		concave points\_worst & 0.7908 \\
		perimeter\_worst & 0.7829 \\
		radius\_worst & 0.7765 \\
		concave points\_mean & 0.7683 \\
		perimeter\_mean & 0.7426 \\
		area\_worst & 0.7338 \\
		radius\_mean & 0.7300 \\
		area\_mean & 0.7090 \\
		concavity\_mean & 0.6863 \\
		concavity\_worst & 0.6492 \\
		compactness\_mean & 0.5965 \\
		compactness\_worst & 0.5910 \\
		radius\_se & 0.5671 \\
		perimeter\_se & 0.5561 \\
		area\_se & 0.5482 \\
		texture\_worst & 0.4569 \\
		smoothness\_worst & 0.4215 \\
		symmetry\_worst & 0.4163 \\
		texture\_mean & 0.4152 \\
		concave points\_se & 0.3907 \\
		smoothness\_mean & 0.3586 \\
		symmetry\_mean & 0.3305 \\
		fractal dimension\_worst & 0.3239 \\
		compactness\_se & 0.2930 \\
		concavity\_se & 0.2378 \\
		fractal dimension\_se & 0.0780 \\
		id & 0.0398 \\
		symmetry\_se & -0.0065 \\
		texture\_se & -0.0083 \\
		fractal dimension\_mean & -0.0128 \\
		smoothness\_se & -0.0670 \\
		\hline
	\end{tabular}
	\caption{Pearson Correlation Coefficients between Features and Diagnosis}
	\label{tab:correlation}
\end{table}

The feature selection process reduced the number of input features from 31 to 10. The classification accuracy before and after feature selection is presented in Table \ref{tab:accuracy_comparison}. The reduced feature set slightly improved the Decision Tree performance from 90.06\% to 90.64\%, indicating that removing weakly correlated features can improve model generalisation.

\begin{table}
	\centering
	\begin{tabular}{l c}
		\hline
		\textbf{Feature Set} & \textbf{Accuracy} \\
		\hline
		Original Features & 90.06\% \\
		Pearson Correlation Feature Selection & 90.64\% \\
		\hline
	\end{tabular}
	\caption{Decision Tree Accuracy Before and After Pearson Feature Selection}
	\label{tab:accuracy_comparison}
\end{table}


\subsubsection{Feature Extraction using PCA}

Principal Component Analysis (PCA) was implemented using the scikit-learn library to reduce the dimensionality of the dataset. Since PCA is sensitive to feature scale, the data was normalised before applying PCA to ensure that features with larger numerical ranges did not dominate the principal components. The same transformation learned from the training data was applied to the testing data to maintain consistency.

The explained variance ratio of each principal component is shown in Table \ref{tab:pca_variance}. The first principal component explains 49.81\% of the variance, while the first ten components together preserve 94.90\% of the total variance. This indicates that PCA can significantly reduce dimensionality while retaining most of the important information from the original dataset.

\begin{table}
	\centering
	\begin{tabular}{c c}
		\hline
		\textbf{Principal Component} & \textbf{Explained Variance Ratio} \\
		\hline
		PC1 & 0.4981 \\
		PC2 & 0.1740 \\
		PC3 & 0.0691 \\
		PC4 & 0.0661 \\
		PC5 & 0.0395 \\
		PC6 & 0.0341 \\
		PC7 & 0.0294 \\
		PC8 & 0.0169 \\
		PC9 & 0.0115 \\
		PC10 & 0.0103 \\
		\hline
	\end{tabular}
	\caption{Explained Variance Ratio of PCA Components}
	\label{tab:pca_variance}
\end{table}

The Decision Tree classifier was evaluated using a fixed hyperparameter setting with a maximum tree depth of 4. The comparison between the original feature space, Pearson correlation feature selection, and PCA based feature extraction is shown in Table \ref{tab:pca_accuracy}. PCA achieved the highest accuracy of 92.40\%, demonstrating that transforming the original features into principal components can improve classification performance.

\begin{table}
	\centering
	\begin{tabular}{l c}
		\hline
		\textbf{Method} & \textbf{Accuracy} \\
		\hline
		Before Feature Selection & 90.06\% \\
		Pearson Correlation Selection & 90.64\% \\
		PCA Feature Extraction & 92.40\% \\
		\hline
	\end{tabular}
	\caption{Accuracy Comparison of Feature Reduction Methods}
	\label{tab:pca_accuracy}
\end{table}

Overall, Pearson correlation feature selection provided a small improvement by removing less relevant features, while PCA achieved a larger improvement by creating new feature combinations that capture the most significant variance in the dataset. The results suggest that PCA was more effective for this dataset because it reduced noise and improved the Decision Tree's ability to separate the two diagnosis classes.

\subsection{Comparison between PCA Feature Extraction and Feature Selection}

To provide a fair comparison between PCA-based feature extraction and Pearson correlation-based feature selection, the number of PCA components was set to match the number of features retained after feature selection. Pearson correlation feature selection reduced the original dataset to 10 features; therefore, PCA was configured to retain 10 principal components. This ensures that both approaches use the same number of dimensions when training the Decision Tree classifier.

A Decision Tree model was trained separately using the PCA-transformed dataset and the feature-selected dataset. The classification accuracy of both approaches was then compared to evaluate the effectiveness of reducing the feature space. The results are shown in Table~\ref{tab:pca_vs_feature_selection}.

\begin{table}[h]
	\centering
	\small
	\begin{tabular}{l c}
		\hline
		\textbf{Dimensionality Reduction Method} & \textbf{Decision Tree Accuracy (\%)} \\
		\hline
		Pearson Correlation (10 features) & 90.64 \\
		PCA (10 components) & 92.40 \\
		\hline
	\end{tabular}
	\caption{Comparison between Feature Selection and PCA using the same number of dimensions}
	\label{tab:pca_vs_feature_selection}
\end{table}

The results show that PCA achieved a higher accuracy than Pearson correlation feature selection, improving the Decision Tree performance from 90.64\% to 92.40\%. This improvement occurs because PCA creates new features by combining information from the original variables, allowing the model to capture hidden patterns and reduce noise within the dataset.

Feature selection has the advantage of maintaining the original feature meanings, making the selected features easier to interpret. For example, the retained features correspond directly to physical measurements such as radius, perimeter, and concavity, which can provide meaningful insights into the diagnosis prediction. Additionally, feature selection reduces computational complexity by removing irrelevant features without transforming the data. However, feature selection may discard useful information because it evaluates features individually and does not consider relationships between multiple features.

In contrast, PCA considers the relationships between all input variables and creates a smaller set of transformed features that preserve the maximum amount of variance. This can improve classification performance and reduce the effect of correlated features. However, PCA components are combinations of the original features and therefore lack direct interpretability. Furthermore, PCA requires feature normalisation before transformation, adding an additional preprocessing step.

Overall, PCA provided better predictive performance for the Decision Tree classifier, while Pearson correlation feature selection provided better interpretability. The choice between these methods depends on the objective of the application: PCA is more suitable when classification accuracy is the priority, whereas feature selection is preferable when understanding the contribution of individual features is important.

\subsection{10-Fold Cross-Validation Evaluation}

To evaluate the reliability and stability of the classification models, 10-fold stratified cross-validation was implemented for both K-Nearest Neighbour (KNN) and Decision Tree classifiers. Stratified cross-validation was used to ensure that each fold maintained a similar distribution of the two diagnosis classes. The KNN classifier was evaluated using two hyperparameter settings, $k=3$ and $k=9$, while the Decision Tree classifier was evaluated using maximum tree depths of 2 and 8.

For each model configuration, the classification accuracy was recorded for every fold. The mean accuracy and standard deviation were calculated to measure overall performance and model stability. A lower standard deviation indicates that the model performance is more consistent across different training and validation splits.

The results of the 10-fold cross-validation are presented in Table~\ref{tab:cross_validation}.

\begin{table*}[h]
	\centering
	\footnotesize
	\begin{tabular}{c c c c c}
		\hline
		\textbf{Fold} & \textbf{KNN (k=3)} & \textbf{KNN (k=9)} & 
		\textbf{Tree (depth=2)} & \textbf{Tree (depth=8)} \\
		\hline
		Fold 1  & 98.25 & 100.00 & 89.47 & 89.47 \\
		Fold 2  & 98.25 & 98.25  & 85.96 & 91.23 \\
		Fold 3  & 98.25 & 96.49  & 89.47 & 92.98 \\
		Fold 4  & 94.74 & 94.74  & 80.70 & 87.72 \\
		Fold 5  & 94.74 & 92.98  & 91.23 & 89.47 \\
		Fold 6  & 92.98 & 92.98  & 94.74 & 98.25 \\
		Fold 7  & 100.00 & 98.25 & 92.98 & 91.23 \\
		Fold 8  & 94.74 & 96.49  & 96.49 & 92.98 \\
		Fold 9  & 96.49 & 98.25  & 92.98 & 92.98 \\
		Fold 10 & 98.21 & 98.21 & 94.64 & 92.86 \\
		\hline
		Mean & 96.66 & 96.66 & 90.87 & 91.92 \\
		Std Dev & 2.14 & 2.28 & 4.49 & 2.74 \\
		\hline
	\end{tabular}
	\caption{10-Fold Cross-Validation Accuracy Results for KNN and Decision Tree Models}
	\label{tab:cross_validation}
\end{table*}

The results show that both KNN configurations achieved the highest classification accuracy, with both $k=3$ and $k=9$ obtaining a mean accuracy of 96.66\%. The similar performance between the two values of $k$ indicates that the choice of neighbourhood size has limited impact on this dataset. However, KNN with $k=3$ achieved a slightly lower standard deviation (2.14\%) compared with $k=9$ (2.28\%), suggesting that the smaller neighbourhood size provides slightly more consistent predictions across different folds.

For the Decision Tree models, increasing the maximum depth from 2 to 8 improved the mean accuracy from 90.87\% to 91.92\%. This improvement occurs because a deeper tree can capture more complex relationships within the dataset. However, the Decision Tree with maximum depth 2 produced the highest standard deviation of 4.49\%, indicating that the shallow tree was more sensitive to changes in the training data. In comparison, the depth-8 Decision Tree achieved a lower standard deviation of 2.74\%, demonstrating improved stability.

Overall, KNN demonstrated the highest performance and reliability among the evaluated models, achieving the best accuracy with relatively low variation between folds. The Decision Tree models provided lower accuracy but offer greater interpretability because the decision-making process can be represented through feature-based rules. The results suggest that KNN is more suitable when predictive performance and stability are prioritised, while Decision Tree models may be preferred when model interpretability is more important.

\section{Statement}

I declare that this assignment is my own work. The implementation, experimentation, model training, and evaluation were completed by me. Artificial intelligence Chatgpt\cite{chatgpt2026} were used only for language assistance during the preparation of this report.

\bibliographystyle{IEEEtran}   % or plain, unsrt, apalike, etc.
\bibliography{references}

\end{document}
